This chapter describes the infrastructure the recommender (Chapters~\ref{ch:behavior}
and~\ref{ch:recommender}) is built and evaluated against: two running Virtual Knowledge Graphs, a
domain-switching development harness, and the recommender's own API contract and persistence
layer. As established in Section~\ref{sec:intro-scope}, none of this chapter's content is itself a
thesis deliverable --- domain and user selection are the responsibility of a separate, external
application. This chapter is included because the recommender genuinely depends on parts of it: the
schema-grounding validation step (Section~\ref{sec:rec-validation}) consumes the machine-readable
schema representations produced here, one of the two evaluation domains (Tourism,
Section~\ref{sec:eval-results}) is only reachable through this infrastructure, and the API contract
described in Section~\ref{sec:arch-api} is the boundary the recommender core is built inside.

\section{Overview}
\label{sec:arch-overview}

Before this work, the NL2SPARQL pipeline queried a single, hardcoded SPARQL endpoint with no
notion of alternate domains, and had no concept of user identity or query persistence. Three layers
were added on top of the unmodified pipeline core: a containerized data/infrastructure layer
(Section~\ref{sec:arch-vkgs}), a domain registry and switching mechanism used only by the internal
development harness (Section~\ref{sec:arch-domain}), and the recommender's own API contract and
query log (Section~\ref{sec:arch-api}), which is the actual thesis-relevant boundary. Because the
NL2SPARQL pipeline's internals --- schema formatting, class/property relevance extraction, entity
linking, SPARQL generation, execution, and result formatting --- are entirely unmodified, any defect
introduced by this work is necessarily located in the new infrastructure, not in query-answering
logic that predates this thesis.

\section{Two Virtual Knowledge Graphs}
\label{sec:arch-vkgs}

Two domains were stood up as isolated Docker Compose stacks, each a relational database container
feeding an Ontop container exposing the data as a VKG over SPARQL (Figure~\ref{fig:arch-phase0}).

\definecolor{dbcolorA}{RGB}{215,235,250}
\definecolor{ontopcolorA}{RGB}{255,235,205}
\definecolor{regcolorA}{RGB}{255,225,225}
\tikzset{blockA/.style={rectangle, draw=colBorder, rounded corners, thick, text width=3.1cm, minimum height=1.1cm, align=center, font=\small}}
\tikzset{smallblockA/.style={rectangle, draw=colBorder, rounded corners, thick, text width=2.5cm, minimum height=0.9cm, align=center, font=\footnotesize}}

\begin{figure}[H]
\centering
\begin{adjustbox}{max width=0.85\textwidth, max totalheight=0.7\textheight, center}
\begin{tikzpicture}[node distance=1.4cm and 2cm]

    \node[blockA, fill=colSource] (pgm) {\code{postgres-movie}\\[2pt]\footnotesize db: \code{netflix\_vkg}\\ port 5433};
    \node[blockA, fill=colOrange, right=of pgm] (om) {\code{ontop-movie}\\[2pt]\footnotesize port 8081};

    \node[blockA, fill=colSource, below=of pgm] (pgt) {\code{postgres-tourism}\\[2pt]\footnotesize db: \code{tourismdb} (PostGIS)\\ port 5434};
    \node[blockA, fill=colOrange, right=of pgt] (ot) {\code{ontop-tourism}\\[2pt]\footnotesize port 8082};

    \node[smallblockA, fill=colPink, below right=1.0cm and -0.3cm of pgt] (jdbc) {\code{docker\slash{}jdbc\slash{}postgresql.jar}\\[2pt]\footnotesize bind-mounted into both Ontop containers};

    \draw[arrow] (pgm) -- node[above,font=\scriptsize]{healthy} (om);
    \draw[arrow] (pgt) -- node[above,font=\scriptsize]{healthy} (ot);
    \draw[arrow, dashed] (jdbc) -- (om);
    \draw[arrow, dashed] (jdbc) -- (ot);

\end{tikzpicture}
\end{adjustbox}
\caption{Docker Compose topology for the two Virtual Knowledge Graphs. Each Ontop container waits
for its database's healthcheck before starting, and both mount the same JDBC driver directory
(dashed arrows) --- a workaround for a missing-driver defect in the base Ontop image.}
\label{fig:arch-phase0}
\end{figure}

\subsection{The Movie domain}

The Movie side uses a synthetic, Faker-generated Netflix-shaped dataset chosen specifically because
it carries per-user behavioral signals (watch history, search logs, reviews) the Tourism domain
lacks --- the substrate Chapter~\ref{ch:behavior} depends on. \code{postgres-movie} is initialized
from six tables (Table~\ref{tab:arch-movie-tables}); no table has a primary key, since the source
dataset contains genuine duplicate id values across every table, a deliberate accommodation rather
than an oversight.

\begin{table}[H]
\centering
\begin{tabular}{lrl}
\toprule
\textbf{Table} & \textbf{Rows} & \textbf{Role} \\
\midrule
\code{movies}              & 1{,}040   & Catalog \\
\code{users}               & 10{,}300  & User master data \\
\code{watch\_history}       & 105{,}000 & Real behavior (implicit feedback) \\
\code{search\_logs}         & 26{,}500  & Stated intent (free-text search) \\
\code{reviews}             & 15{,}450  & Explicit stated opinions \\
\code{recommendation\_logs} & 52{,}000  & System-suggested vs.\ consumed \\
\bottomrule
\end{tabular}
\caption{Movie-domain tables, with row counts verified exactly against the source CSVs.}
\label{tab:arch-movie-tables}
\end{table}

The pre-existing OBDA mapping for this domain required two categories of fixes before it was usable:
a bug fix (a mapping selected a non-existent column, \code{watch\_duration}, where the real column
is \code{watch\_duration\_minutes}), and a coverage extension --- to 17 mappings in total --- adding
the columns a later phase's search-to-watch temporal join needs (\code{search\_date},
\code{watch\_date}, \code{action}, \code{device\_type}), plus a new object property connecting
\code{WatchSession} to \code{User} that had not previously existed despite the equivalent
\code{WatchSession}$\to$\code{Movie} relationship being present.

\subsection{The Tourism domain}

The Tourism side reuses a real, 293\,MB production PostgreSQL dump of Verona's tourism
application, together with an ontology that --- unlike the Movie domain --- already embeds SHACL
node shapes per class, a direct structural match for the NL2SPARQL pipeline's schema loader
(Section~\ref{sec:arch-domain}). \code{postgres-tourism} runs on a PostGIS-enabled image, since the
source database uses geometry/geography columns a plain Postgres image cannot represent.

\begin{table}[H]
\centering
\begin{tabular}{lr}
\toprule
\textbf{Table} & \textbf{Rows} \\
\midrule
\code{art}      & 41 \\
\code{event}    & 1{,}416 \\
\code{calendar} & 10{,}640 \\
\code{location} & 1{,}600 \\
\bottomrule
\end{tabular}
\caption{Tourism-domain tables, row counts verified exactly against the source backup.}
\label{tab:arch-tourism-tables}
\end{table}

Two mappings in the original OBDA file referenced a table, \code{art\_images}, that does not exist
anywhere in the backup (only a differently-shaped legacy table exists). Rather than reconstructing
the missing table, these two mappings were dropped --- a decision made explicitly with the
supervisor's sign-off, since art images are not needed for this thesis's demonstration. The
resulting classes remain declared in the ontology for documentation purposes but are unpopulated.

\paragraph{An explicit scope decision.} The backup also contains a large behavioral table,
\code{log\_vc} (4.4 million rows of VeronaCard tap-in visits), which could in principle serve as a
Tourism-domain behavioral signal analogous to the Movie domain's watch history. However, it
identifies an anonymous card, not a named person, and the database's only user-like table has just
4 administrative rows, not tourists. It was decided that Tourism stays catalog-query-only for this
thesis: \code{log\_vc} is not mapped into the ontology, and all user-selection/behavior-aware work
in later chapters is carried entirely by the Movie domain. This is a load-bearing scope decision,
not an oversight --- it bounds what Chapters~\ref{ch:behavior},~\ref{ch:recommender},
and~\ref{ch:eval} can demonstrate on the Tourism side, and is referenced repeatedly in those
chapters as the reason certain results are Movie-only.

\section{Domain-Switching Infrastructure}
\label{sec:arch-domain}

A declarative KG registry (\code{kg\_config.yaml}) lists the available knowledge graphs by name,
schema-file path, SPARQL endpoint, and an \code{enabled} flag, read by a small loader function that
returns only enabled entries. This registry, and the internal development harness that consumes it
to let a developer manually switch domains while exercising the recommender, are not thesis
deliverables in themselves (Section~\ref{sec:intro-scope}) and are not documented further here.

One infrastructural fix in this layer is worth recording because it directly affects
Section~\ref{sec:rec-validation}'s later results: the NL2SPARQL pipeline's schema loader builds its
machine-readable schema representation (``m-schema'') entirely from SHACL \code{sh:NodeShape}
triples, with no fallback for plain \code{rdfs:domain}/\code{rdfs:range} declarations. The Movie
ontology, unlike the Tourism ontology, originally had no SHACL shapes at all, so feeding it directly
to the parser silently produced an empty schema (Figure~\ref{fig:arch-shacl}). A one-time conversion
script generates the missing shapes mechanically from the ontology's existing \code{rdfs:domain}
declarations, producing \code{movie\_mschema.json} (11 classes) alongside the Tourism domain's
already-SHACL-native \code{tourism\_mschema.json} (7 classes). Both files are the direct input to
Section~\ref{sec:rec-validation}'s schema-grounding check.

\definecolor{ownercolorB}{RGB}{215,235,250}
\begin{figure}[H]
\centering
\begin{adjustbox}{max width=0.85\textwidth, max totalheight=0.6\textheight, center}
\begin{tikzpicture}[node distance=1.3cm and 1.8cm]
    \node[block, fill=colPurple] (owl) {\code{netflix\_ontology.owl}\\[2pt]\footnotesize plain OWL, rdfs:domain/range only};
    \node[block, fill=colPink, right=of owl] (script) {\code{generate\_shacl\_shapes.py}\\[2pt]\footnotesize one-time conversion script};
    \node[block, fill=colPurple, right=of script] (shapes) {\code{netflix\_shapes.ttl}\\[2pt]\footnotesize 11 sh:NodeShapes};
    \node[block, fill=colGreen, below=of script, text width=4.2cm, font=\footnotesize] (parser) {SHACL schema parser};
    \node[block, fill=colGreen, right=of parser] (mschema) {\code{movie\_mschema.json}\\[2pt]\footnotesize 11 classes};

    \draw[arrow] (owl) -- (script);
    \draw[arrow] (script) -- (shapes);
    \draw[arrow] (shapes) -- (parser);
    \draw[arrow] (parser) -- (mschema);
\end{tikzpicture}
\end{adjustbox}
\caption{The Movie ontology has no native SHACL shapes, so a one-time script generates them from
its existing \code{rdfs:domain}/\code{rdfs:range} declarations before the schema parser can consume
it. This m-schema JSON is the same artifact Section~\ref{sec:rec-validation}'s schema-grounding
check embeds and thresholds against.}
\label{fig:arch-shacl}
\end{figure}

\section{The Recommender API Contract and Query Log}
\label{sec:arch-api}

This is the layer the recommender core (Chapter~\ref{ch:recommender}) is built inside, and the
only genuinely thesis-relevant infrastructure in this chapter beyond the two VKGs above. Before any
recommendation logic exists, two things are needed: a fixed request/response contract external
callers can integrate against, and a place to persist queries, since the recommender's core
technique depends on retrieving a user's own past queries.

\subsection{The contract}

\begin{lstlisting}[language=Python, caption={The recommender's request/response contract. One
field, \code{interaction\_data}, is deliberately left untyped: its real shape depends on what the
external application actually logs per user, a dependency that was never confirmed during this
thesis (Section~\ref{sec:conclusion-limits}).}]
class RecommendRequest(BaseModel):
    user_id: str
    domain: str
    query_text: str
    generated_sparql: Optional[str] = None
    result_summary: Optional[Any] = None
    interaction_data: Optional[dict] = None

class RecommendResponse(BaseModel):
    suggestions: list[str] = Field(default_factory=list)
    query_log_id: Optional[int] = None
\end{lstlisting}

\subsection{The query log}

A single SQLite table, indexed on \code{(user\_id, domain)}, appended to on every request. SQLite
was chosen deliberately over a heavier database: a single-file, zero-infrastructure store is
appropriate for a thesis-scale log that does not need to be shared across processes beyond one
local API server, consistent with this project's broader pattern of using the lightest tool that
fits. Two functions cover the read and write side: \code{log\_query()} appends one row (with
automatic JSON-serialization of structured result summaries), and \code{get\_user\_history()}
retrieves the most recent rows for a user, optionally scoped to one domain --- the function
Section~\ref{sec:rec-retrieval}'s retrieval step calls to build its candidate pool.

\subsection{One shared entry point}

Both the real HTTP endpoint (\code{POST /recommend}, exposed via FastAPI, with OpenAPI
documentation auto-served at \code{/docs}) and the internal development harness call a single
function, \code{handle\_recommend()} (Figure~\ref{fig:arch-api}). This was a deliberate design
choice, not a literal reading of the original plan (which could be read as requiring the harness to
make a real HTTP call to a separately-running server): the request/response \emph{shape} being
validated is identical either way, since both paths go through the same pydantic models and the
same handler logic, and requiring a developer to keep a second server process running just to see a
suggestions panel render is friction the contract does not need to impose. External callers are
unaffected by this choice --- they only ever see the HTTP surface.

\definecolor{apicolorC}{RGB}{255,235,205}
\begin{figure}[H]
\centering
\begin{adjustbox}{max width=\textwidth, max totalheight=0.85\textheight, center}
\begin{tikzpicture}[node distance=1.0cm and 1.4cm]

    \node[block, fill=colOrange] (http) {External caller\\[2pt]\footnotesize \code{POST /recommend} over HTTP};
    \node[block, fill=colGreen, right=1.8cm of http] (gui) {Development harness\\[2pt]\footnotesize in-process call};

    \node[smallblock, fill=colOrange, below=of http] (app) {FastAPI app\\[2pt]\footnotesize \code{/recommend}, \code{/health}};

    \node[block, fill=colPurple, below=1.6cm of http, xshift=2.5cm] (handle) {\code{handle\_recommend()}\\[2pt]\footnotesize shared entry point};

    \node[block, fill=colPink, below=1.1cm of handle] (log) {\code{log\_query()}\\[2pt]\footnotesize $\to$ SQLite};

    \node[smallblock, fill=colPink, below=of log] (db) {\code{query\_log.db}\\[2pt]\footnotesize indexed on (user\_id, domain)};

    \node[block, fill=colPurple, right=1.6cm of log] (resp) {\code{RecommendResponse}\\[2pt]\footnotesize the recommender core (Ch.~\ref{ch:recommender}) fills this in};

    \draw[arrow] (http) -- (app);
    \draw[arrow] (app) -- (handle);
    \draw[arrow] (gui) -- (handle);
    \draw[arrow] (handle) -- (log);
    \draw[arrow] (log) -- (db);
    \draw[arrow] (handle) -- (resp);

\end{tikzpicture}
\end{adjustbox}
\caption{Both the real HTTP endpoint and the internal development harness converge on
\code{handle\_recommend()}, which validates the request, persists it, and returns a response. This
is the extension point the recommender core (Chapter~\ref{ch:recommender}) is built inside, without
any further change to the contract, the log store, or the harness.}
\label{fig:arch-api}
\end{figure}

\subsection{Live rating capture}
\label{sec:arch-ratings}

The offline ablation in Chapter~\ref{ch:eval} scores suggestions against held-out historical
behavior, not against feedback from a real user who actually saw them. To make that second,
complementary form of evaluation possible, a second endpoint, \code{POST /rate}, was added
alongside \code{POST /recommend}:

\begin{lstlisting}[language=Python, caption={The rating contract. \code{query\_log\_id} is the id
returned by the \code{/recommend} call that produced the rated suggestion, so a rating can be
joined back to the request it responds to.}]
class RateRequest(BaseModel):
    user_id: str
    domain: str
    suggestion: str
    stars: int = Field(ge=1, le=5)
    query_log_id: Optional[int] = None

class RateResponse(BaseModel):
    rating_id: int
\end{lstlisting}

\code{handle\_rate()} is deliberately the simplest handler in the API: unlike
\code{handle\_recommend()}, it orchestrates no pipeline --- it validates the 1--5 star rating
(Pydantic-enforced) and appends one row to a second SQLite store, \code{ratings.db}, mirroring the
query log's single-file, zero-infrastructure pattern (\code{ratings.py}, indexed on \code{domain}
and \code{query\_log\_id}). A companion script, \code{ratings\_summary.py}, reads that table back
and reports mean stars and rating counts overall, per domain, and per suggestion. Both the endpoint
and the summary script are domain-agnostic by construction --- neither depends on a behavioral
profile existing for the rated user --- so the mechanism already covers Tourism as well as Movie,
unlike the ablation in Chapter~\ref{ch:eval}. No real ratings had been collected by the time this
thesis was written, so this section documents a working, exercised mechanism (demonstrated
end-to-end by \code{call\_recommender.py}: call \code{/recommend}, then rate its first suggestion),
not a results table; using it to gather real user feedback is future work
(Section~\ref{sec:conclusion-future}).

\section{Summary}

This chapter described the infrastructure the recommender depends on but does not itself
contribute: two running VKGs (Movie, behaviorally rich but synthetic; Tourism, real but
catalog-only by explicit scope decision), a domain-switching development harness, and --- the one
piece that is genuinely load-bearing for the thesis --- a fixed API contract, persistent query log,
and live rating store that Chapter~\ref{ch:recommender}'s recommender core is built entirely
inside, with no further change to this boundary.
